% NSF GRFP Research Statement
% Formatting requirements: 1-inch margins, 11pt or 12pt font, single-spaced
% Maximum 2 pages for Graduate Research Plan Statement

\documentclass[11pt]{article}

% Page formatting - NSF GRFP requires 1-inch margins
\usepackage[margin=1in]{geometry}
\usepackage{times} % Times New Roman font (similar to NSF standard)
\usepackage{setspace}
% \usepackage{enumitem}
% \usepackage{hyperref}

% Remove page numbers if desired (check current NSF requirements)
\pagestyle{plain}

% Reduce spacing
\setlength{\parindent}{0pt}
\setlength{\parskip}{6pt}

% Section formatting
% \usepackage{titlesec}
% \titleformat{\section}{\normalfont\fontsize{12}{14}\bfseries}{\thesection}{1em}{}
% \titleformat{\subsection}{\normalfont\fontsize{11}{13}\bfseries}{\thesubsection}{1em}{}
% \titlespacing*{\section}{0pt}{8pt}{4pt}
% \titlespacing*{\subsection}{0pt}{6pt}{2pt}
% 
\begin{document}

\begin{center}
\textbf{\large Graduate Research Plan Statement}\\
\textbf{Provably Reliable and Interpretable AI Systems for Therapeutic Enzyme Engineering}
\end{center}

\section*{My Research Vision} 

Publicly funded biological data—built on billions in taxpayer investment through
NSF, NIH, and other agencies—represents one of science’s greatest untapped
assets. Yet despite its scale and richness, we lack AI systems capable of {\bf
reasoning over this data with rigor, transparency, and reliability}. Today’s
dominant models are black boxes: they predict, but don’t explain; they
correlate, but don’t guarantee. In therapeutic contexts—where a single error
could mean toxic side effects or failed treatments—this opacity is unacceptable.
My PhD research will pioneer the first generation of {\bf Verifiable,
Interpretable, and Provable AI agents (VIP-AI)} for therapeutic enzyme
engineering, closing the critical trust gap between AI and clinical application.

I will integrate cutting-edge neural architectures—protein language models,
state-space models (e.g., Mamba), and biophysically constrained attention
mechanisms—with formal verification and mathematical priors from structural
biology to build models that deliver not just {\bf accuracy}, but {\bf certified safety
bounds} and {\bf human-understandable mechanistic explanations}. This work advances
core computer science in trustworthy AI while directly enabling safer, more
rational design of {\bf enzyme therapeutics—a rapidly growing frontier in precision
medicine for cancer, aging, and genetic disease}. By transforming public data
into provably reliable knowledge, my research ensures AI earns its place in the
lab and the clinic: not as a mysterious oracle, but as a transparent,
accountable partner in human health.




\section*{Intellectual Merit}

\subsection*{Goal 1: Verifiable Neural Architectures for Mutation Effect Prediction}

\textbf{Research Question:} How can we design neural architectures with formal
verification guarantees that accurately predict mutation effects while
respecting fundamental biochemical constraints?

\textbf{Approach:} Current approaches to predicting mutation effects ($\Delta
\Delta G$ stability changes, catalytic efficiency) rely on transformers (ESM-2,
ProteinMPNN) or physics-based methods (Rosetta, FoldX), but lack formal
reliability guarantees and ignore physical constraints. I will develop a novel
hybrid architecture combining: (1) \textit{Mamba-based sequence encoders}
leveraging linear-time complexity to efficiently process long protein sequences
(500+ residues) while maintaining biophysically meaningful hidden states, (2)
\textit{constrained attention mechanisms} where QK-V attention weights respect
biochemical principles through differentiable projection operators from convex
optimization, and (3) \textit{formal verification layers} that provide certified
uncertainty bounds using abstract interpretation and SMT solvers, outputting
provable guarantees such as "With 95\% confidence, the actual $\Delta \Delta G$ lies within
[-0.5, 0.3] kcal/mol."

\textbf{Expected Outcomes:} This will produce theoretical guarantees on model
reliability quantified through PAC-learning bounds adapted to structured
biochemical data, empirically validated architectures on benchmarks (SKEMPI 2.0,
Mega-scale), and open-source verification tools for the computational biology
community.

\subsection*{Goal 2: Interpretable Reasoning for Enzyme Mechanism Understanding}

\textbf{Research Question:} Can we develop AI systems that generate verifiable,
mechanistically accurate explanations that structural biologists can validate?

\textbf{Approach:} Understanding \textit{why} mutations improve function is
crucial for rational design. I will create neuro-symbolic architectures
performing explicit reasoning over enzyme mechanisms through: (1)
\textit{structured reasoning with graph neural networks} representing
biochemical mechanisms as computational graphs where nodes denote active sites,
transition states, and cofactors, with message-passing generating "chains of
thought" mirroring expert reasoning, (2) \textit{attention-based mechanism
discovery} using mechanistic interpretability to identify which attention heads
in protein language models correspond to specific biochemical features
(catalytic triads, electrostatic complementarity), and (3) \textit{theorem
prover integration} using differentiable theorem provers to verify explanations
against biochemical rules and quantum mechanical calculations.

\textbf{Expected Outcomes:} This advances AI interpretability theory while
providing practical tools, including formalization of enzyme mechanisms as
logical constraints, models generating verifiable explanations with $>$80\%
expert agreement, and case studies designing improved therapeutic enzymes (e.g.,
L-asparaginase for acute lymphoblastic leukemia).

\subsection*{Goal 3: Certified Safety Guarantees for Therapeutic Enzyme Design}

\textbf{Research Question:} How can we provide formal safety certificates for
AI-designed therapeutic enzymes, ensuring substrate specificity, stability, and
lack of off-target activity?

\textbf{Approach:} Therapeutic enzymes must satisfy stringent safety
criteria—enzymes for prodrug activation must be exquisitely substrate-specific,
while those degrading toxic metabolites must not cleave essential proteins.
Current generative models (ProteinMPNN, ESM-IF) offer no such guarantees. I will
develop: (1) \textit{constraint-guided generative models} incorporating hard
constraints during sequence generation using projected gradient descent and
Lagrangian optimization, provably satisfying constraints like "$>$100-fold
binding selectivity," (2) \textit{uncertainty quantification with certified
coverage} through conformal prediction methods providing guaranteed coverage
("With 99\% probability, true activity falls within the predicted set"), and (3)
\textit{multi-objective optimization with Pareto guarantees} identifying
provably Pareto-optimal enzyme variants balancing efficiency, stability, and
specificity using Bayesian optimization.

\textbf{Expected Outcomes:} This establishes theoretical foundations for
trustworthy AI in protein engineering: formal definitions of safety expressed in
temporal logic, algorithms with provable constraint satisfaction guarantees, and
experimental validation demonstrating AI-certified enzymes meet predicted
specifications.

\section*{Broader Impacts}

\textbf{Advancing Human Health and Longevity:} This research directly addresses
aging and disease challenges. Enzyme replacement therapies for lysosomal storage
disorders cost $>$\$200,000 per patient annually; improved variants could reduce
costs and improve efficacy. My focus on longevity-relevant enzymes (DNA repair
glycosylases, NAD+ biosynthesis enzymes, senolytic proteases) accelerates paths
from computational prediction to clinical application.

\textbf{Democratizing Therapeutic Development:} Current enzyme engineering
requires expensive high-throughput screening accessible only to well-funded
institutions. Open-source, reliable AI tools democratize this capability,
enabling researchers at resource-limited institutions and developing countries
to design therapeutic enzymes computationally. I will release all software with
comprehensive documentation and tutorials for structural biologists without
extensive ML backgrounds.

\textbf{Interdisciplinary Education and Mentorship:} I am committed to training
researchers at the intersection of AI and biology through: (1) workshop series
"Trustworthy AI for Drug Discovery" bringing together ML researchers and
biochemists, (2) mentoring undergraduates from underrepresented groups through
university STEM diversity initiatives, and (3) curriculum development
integrating formal verification into computational biology courses.

\textbf{Responsible AI Development:} By prioritizing interpretability and
verification, this work contributes to AI safety and responsible innovation.
Methods developed—certified predictions, mechanistic explanations, formal
guarantees—transfer to other high-stakes medical AI applications. I will engage
with bioethics communities to ensure appropriate safeguards and equitable
access.

\section*{Conclusion}

The convergence of efficient neural architectures, formal verification methods,
and computational biology creates an unprecedented opportunity to transform
therapeutic enzyme design. My research advances fundamental computer
science—developing theory for verifiable learning, interpretable AI, and
constrained optimization on structured spaces—while addressing human health
challenges. By creating AI systems that are provably reliable and
understandable, I aim to establish a new paradigm for trustworthy machine
learning in life sciences, embodying NSF's mission to support research advancing
both scientific knowledge and societal well-being.

\end{document}

% Enzyme therapeutics represent a rapidly growing
% frontier in precision medicine, with applications ranging from cancer treatment
% to age-related disease intervention. However, rational enzyme design faces a
% fundamental challenge: current machine learning models lack the reliability
% guaranties and interpretability needed for therapeutic applications. When
% prediction errors could mean the difference between life-saving treatments and
% toxic therapies, we cannot rely on black-box models that offer no formal
% guaranties or mechanistic understanding.
% 
% My research aims to develop the first generation of \textbf{V.I.P.}
% --Verifiable, Interpretable, and Provable  AI systems for therapeutic enzyme
% engineering. By integrating advances in efficient neural architectures—including
% state-space models (Mamba), novel attention mechanisms, and protein language
% models—with formal verification methods and mathematical constraints from
% biophysics, I will create AI systems that provide certified safety guarantees
% and human-interpretable explanations. This work advances fundamental computer
% science theory while addressing pressing societal challenges in human health and
% longevity.

Every year, U.S. taxpayers invest billions in scientific discovery—fueling
breakthroughs in genomics, structural biology, and enzyme engineering through
visionary agencies like the NSF and NIH. The result? A goldmine of public
biological data, richer and more detailed than ever before.

Yet here’s the paradox: while we’ve mastered generating data, we’re still
fumbling in the dark when it comes to reasoning with it. Most AI tools used in
biomedicine today are black boxes—impressive at pattern-matching, but silent on
why they make a prediction or whether we can trust it. In drug design, that
silence isn’t just inconvenient—it can be dangerous.

That’s why my PhD research will build the first generation of V.I.P. AI:
Verifiable, Interpretable, and Provable systems engineered from the ground up
for therapeutic enzyme design. Imagine an AI that doesn’t just say “this enzyme
will work”—but shows you how, proves it won’t misfire, and explains its logic in
terms a biologist (or clinician) can understand.

I’ll fuse the best of modern deep learning—protein language models, state-space
architectures like Mamba, and biophysically grounded attention—with the rigor of
formal verification and mathematical constraints from molecular physics. The
result? AI that delivers not just predictions, but certified guarantees: “This
design is safe within these bounds,” or “This mutation disrupts catalysis—here’s
the evidence.”

This isn’t just incremental AI—it’s a new paradigm for trustworthy science.
Enzyme therapeutics are already revolutionizing precision medicine, from
targeting cancer cells to reversing age-related decline. But to unlock their
full potential, we need models we can trust like a colleague, not just use like
a calculator.

My vision is clear: AI that earns its place in the lab—and in the clinic—by
being as transparent as it is intelligent. By turning public data into provably
reliable knowledge, this work will advance the frontiers of computer science
while directly serving human health, longevity, and the public good.


formal tone:

Publicly funded biological data—generated through billions of dollars in
taxpayer investment via agencies like the NSF and NIH—represents one of the most
valuable scientific resources of our time. Yet surprisingly little research has
focused on developing AI systems capable of reasoning over this data with the
rigor, transparency, and reliability demanded by high-stakes applications in
human health.

My PhD research will pioneer the first generation of V.I.P. AI
systems—Verifiable, Interpretable, and Provable—specifically designed for
therapeutic enzyme engineering. By unifying cutting-edge neural architectures
(including state-space models like Mamba, protein language models, and
biophysically informed attention mechanisms) with formal verification techniques
and mathematical constraints from structural biology, I will build AI that not
only predicts enzyme function with high accuracy but also provides certified
safety guarantees and human-understandable explanations of its reasoning.

This work bridges fundamental computer science—advancing theories in trustworthy
AI, formal methods, and geometric deep learning—with urgent societal needs in
precision medicine. Enzyme therapeutics are emerging as a transformative
frontier, with applications from oncology to age-related disease intervention.
Yet rational enzyme design remains bottlenecked by a critical gap: current
machine learning models are black boxes that offer neither reliability
guarantees nor mechanistic insight. In therapeutic contexts—where a prediction
error could mean the difference between a life-saving treatment and a toxic
outcome—trust must be earned through verifiability, not assumed through accuracy
alone. My research aims to close this gap, ensuring that AI becomes a truly
reliable partner in the quest for safer, more effective medicines.


